\subsection{Limitations}\label{sec:limits}
All quantitative results come from simulation. LunarSim reproduces the statistical character of
lunar relief and the dominant photometric effects, but it omits several things: resolved near-field rocks below the $0.3$\,m DEM spacing (\cref{fig:apollo}), multiple
scattering beyond a constant ambient term, rover self-shadowing and dust, lens flare, exposure
bracketing, and the full range of geological units. Absolute recall values should therefore not be
read as predictions of on-Moon performance. The \emph{relative} comparisons are made under
identical conditions and with paired queries. Because the backbone is frozen and no method sees
test scenes, we expect these comparisons to carry over better than the absolute values, but this
must be confirmed on real imagery. The backbone was frozen because of compute constraints, and
joint fine-tuning may change the ranking between heads. The FGW structural term uses image-plane
distances, which are only approximately preserved under large viewpoint changes. Lifting database
keypoints to 3-D with DEM depth is a natural extension.

% ======================================================================
\section{Conclusion}\label{sec:conclusion}
% ======================================================================
Galatea treats both stages of lunar place recognition as optimal-transport problems. For
aggregation, a semi-relaxed unbalanced formulation unifies the NetVLAD and SALAD assignments under
one parameter. It comes with an explicit convergence rate, a closed-form view of the dust-bin, and
an invariance result showing that intra-normalised descriptors see the plan only through token
weights. For re-ranking, partial fused Gromov--Wasserstein matching adds the relative layout of
keypoints to their appearance, and a saturating fusion merges local and global evidence. On the
LunarSim-VPR benchmark, three findings are useful for practitioners:
(i)~terrestrial models, including foundation-model ones, lose more than half their recall under
large Sun changes;
(ii)~illumination-diverse training data is worth far more than the choice of aggregation layer,
and, as predicted by the theory, balanced and unbalanced assignments perform alike;
(iii)~fusing local with global evidence adds five to seven points of Recall@1 regardless of the
local matcher, whereas re-ranking by local evidence alone helps little.
The next steps are evaluation on DEM-registered real imagery, joint fine-tuning of the backbone,
and 3-D structural terms that use the database's DEM depth.
